\documentclass[11pt,a4paper]{article}
% Preambulo comun de los reportes de Tecnologias Celulares.
% Cada main.tex hace: \documentclass ... % Preambulo comun de los reportes de Tecnologias Celulares.
% Cada main.tex hace: \documentclass ... % Preambulo comun de los reportes de Tecnologias Celulares.
% Cada main.tex hace: \documentclass ... \input{../comun/preambulo.tex}
\usepackage[utf8]{inputenc}
\usepackage[T1]{fontenc}
\usepackage[spanish,mexico]{babel}
\usepackage{graphicx}
\usepackage[margin=2.5cm]{geometry}
\usepackage{amsmath}
\usepackage{amssymb}
\usepackage{booktabs}
\usepackage{tabularx}
\usepackage{array}
\usepackage{float}
\usepackage{setspace}
\usepackage{parskip}
\usepackage{fancyhdr}
\usepackage{textcomp}
\usepackage{xcolor}
\usepackage{caption}
\usepackage{fvextra}
\usepackage[hidelinks]{hyperref}

\onehalfspacing
\setlength{\parskip}{6pt}
\sloppy
\setlength{\emergencystretch}{3em}
\setlength{\headheight}{40pt}
\captionsetup{font=small,labelfont=bf}
\setlength{\tabcolsep}{6pt}

\pagestyle{fancy}
\fancyhf{}
\fancyhead[L]{\includegraphics[width=1cm]{/home/yovick/repos/UAZ-IEI/Assets/UIE.png}}
\fancyhead[R]{\small Tecnologías Celulares}
\fancyfoot[C]{\thepage}
\renewcommand{\headrulewidth}{0.4pt}

\DefineVerbatimEnvironment{codigo}{Verbatim}{%
  breaklines=true,breakanywhere=true,fontsize=\small,frame=single,
  framesep=3pt,rulecolor=\color{gray!50}}

% Figura centrada con pie: \figura[ancho]{archivo}{leyenda}
\newcommand{\figura}[3][0.85]{%
  \begin{center}
    \includegraphics[width=#1\linewidth]{#2}\\[2pt]
    {\small #3}
  \end{center}}

% Portada: \portada{numero}{titulo de la practica}
\newcommand{\portada}[2]{%
\begin{titlepage}
  \thispagestyle{empty}
  \begin{center}
    \vspace*{1cm}
    \includegraphics[width=4.5cm]{/home/yovick/repos/UAZ-IEI/Assets/uaz.png}
    \vspace{0.8cm}

    {\large Universidad Autónoma de Zacatecas}\\[4pt]
    {\large Unidad Académica de Ingeniería Eléctrica}\\[4pt]
    {\normalsize Ingeniería Electrónica Industrial con especialidad en Telecomunicaciones}

    \vspace{1.8cm}
    {\LARGE\bfseries Práctica #1}\\[10pt]
    {\Large\bfseries #2}
  \end{center}

  \vspace{2.2cm}
  \begin{center}
    \begin{tabular}{rl}
      \textbf{Materia:} & Tecnologías Celulares\\
      \textbf{Docente:} & Dr. Víktor Iván Rodríguez Abdalá\\
      \textbf{Alumno:}  & Arnold Yovick Ríos Zaldívar\\
      \textbf{Fecha:}   & Noviembre de 2026\\
    \end{tabular}
  \end{center}
  \vfill
\end{titlepage}}

\usepackage[utf8]{inputenc}
\usepackage[T1]{fontenc}
\usepackage[spanish,mexico]{babel}
\usepackage{graphicx}
\usepackage[margin=2.5cm]{geometry}
\usepackage{amsmath}
\usepackage{amssymb}
\usepackage{booktabs}
\usepackage{tabularx}
\usepackage{array}
\usepackage{float}
\usepackage{setspace}
\usepackage{parskip}
\usepackage{fancyhdr}
\usepackage{textcomp}
\usepackage{xcolor}
\usepackage{caption}
\usepackage{fvextra}
\usepackage[hidelinks]{hyperref}

\onehalfspacing
\setlength{\parskip}{6pt}
\sloppy
\setlength{\emergencystretch}{3em}
\setlength{\headheight}{40pt}
\captionsetup{font=small,labelfont=bf}
\setlength{\tabcolsep}{6pt}

\pagestyle{fancy}
\fancyhf{}
\fancyhead[L]{\includegraphics[width=1cm]{/home/yovick/repos/UAZ-IEI/Assets/UIE.png}}
\fancyhead[R]{\small Tecnologías Celulares}
\fancyfoot[C]{\thepage}
\renewcommand{\headrulewidth}{0.4pt}

\DefineVerbatimEnvironment{codigo}{Verbatim}{%
  breaklines=true,breakanywhere=true,fontsize=\small,frame=single,
  framesep=3pt,rulecolor=\color{gray!50}}

% Figura centrada con pie: \figura[ancho]{archivo}{leyenda}
\newcommand{\figura}[3][0.85]{%
  \begin{center}
    \includegraphics[width=#1\linewidth]{#2}\\[2pt]
    {\small #3}
  \end{center}}

% Portada: \portada{numero}{titulo de la practica}
\newcommand{\portada}[2]{%
\begin{titlepage}
  \thispagestyle{empty}
  \begin{center}
    \vspace*{1cm}
    \includegraphics[width=4.5cm]{/home/yovick/repos/UAZ-IEI/Assets/uaz.png}
    \vspace{0.8cm}

    {\large Universidad Autónoma de Zacatecas}\\[4pt]
    {\large Unidad Académica de Ingeniería Eléctrica}\\[4pt]
    {\normalsize Ingeniería Electrónica Industrial con especialidad en Telecomunicaciones}

    \vspace{1.8cm}
    {\LARGE\bfseries Práctica #1}\\[10pt]
    {\Large\bfseries #2}
  \end{center}

  \vspace{2.2cm}
  \begin{center}
    \begin{tabular}{rl}
      \textbf{Materia:} & Tecnologías Celulares\\
      \textbf{Docente:} & Dr. Víktor Iván Rodríguez Abdalá\\
      \textbf{Alumno:}  & Arnold Yovick Ríos Zaldívar\\
      \textbf{Fecha:}   & Noviembre de 2026\\
    \end{tabular}
  \end{center}
  \vfill
\end{titlepage}}

\usepackage[utf8]{inputenc}
\usepackage[T1]{fontenc}
\usepackage[spanish,mexico]{babel}
\usepackage{graphicx}
\usepackage[margin=2.5cm]{geometry}
\usepackage{amsmath}
\usepackage{amssymb}
\usepackage{booktabs}
\usepackage{tabularx}
\usepackage{array}
\usepackage{float}
\usepackage{setspace}
\usepackage{parskip}
\usepackage{fancyhdr}
\usepackage{textcomp}
\usepackage{xcolor}
\usepackage{caption}
\usepackage{fvextra}
\usepackage[hidelinks]{hyperref}

\onehalfspacing
\setlength{\parskip}{6pt}
\sloppy
\setlength{\emergencystretch}{3em}
\setlength{\headheight}{40pt}
\captionsetup{font=small,labelfont=bf}
\setlength{\tabcolsep}{6pt}

\pagestyle{fancy}
\fancyhf{}
\fancyhead[L]{\includegraphics[width=1cm]{/home/yovick/repos/UAZ-IEI/Assets/UIE.png}}
\fancyhead[R]{\small Tecnologías Celulares}
\fancyfoot[C]{\thepage}
\renewcommand{\headrulewidth}{0.4pt}

\DefineVerbatimEnvironment{codigo}{Verbatim}{%
  breaklines=true,breakanywhere=true,fontsize=\small,frame=single,
  framesep=3pt,rulecolor=\color{gray!50}}

% Figura centrada con pie: \figura[ancho]{archivo}{leyenda}
\newcommand{\figura}[3][0.85]{%
  \begin{center}
    \includegraphics[width=#1\linewidth]{#2}\\[2pt]
    {\small #3}
  \end{center}}

% Portada: \portada{numero}{titulo de la practica}
\newcommand{\portada}[2]{%
\begin{titlepage}
  \thispagestyle{empty}
  \begin{center}
    \vspace*{1cm}
    \includegraphics[width=4.5cm]{/home/yovick/repos/UAZ-IEI/Assets/uaz.png}
    \vspace{0.8cm}

    {\large Universidad Autónoma de Zacatecas}\\[4pt]
    {\large Unidad Académica de Ingeniería Eléctrica}\\[4pt]
    {\normalsize Ingeniería Electrónica Industrial con especialidad en Telecomunicaciones}

    \vspace{1.8cm}
    {\LARGE\bfseries Práctica #1}\\[10pt]
    {\Large\bfseries #2}
  \end{center}

  \vspace{2.2cm}
  \begin{center}
    \begin{tabular}{rl}
      \textbf{Materia:} & Tecnologías Celulares\\
      \textbf{Docente:} & Dr. Víktor Iván Rodríguez Abdalá\\
      \textbf{Alumno:}  & Arnold Yovick Ríos Zaldívar\\
      \textbf{Fecha:}   & Noviembre de 2026\\
    \end{tabular}
  \end{center}
  \vfill
\end{titlepage}}


\begin{document}

\portada{9}{Modelado de canales MIMO y capacidad de sistemas de antena múltiple}

\textbf{Script:} \texttt{practica9.py} (Python con NumPy)

\section{Objetivo}

Modelar canales MIMO Rayleigh, generar matrices de canal aleatorias y evaluar la capacidad
de canal para configuraciones SISO, 2x2 y 4x4, analizando la influencia del número de
antenas y de la relación señal a ruido.

\section{Fundamento teórico}

Un sistema MIMO se describe por
\begin{equation}
  \mathbf{y} = \mathbf{H}\mathbf{x} + \mathbf{n},
\end{equation}
donde $\mathbf{H}\in\mathbb{C}^{N_r\times N_t}$ es la matriz de canal, $\mathbf{x}$ el
vector transmitido y $\mathbf{n}$ el ruido gaussiano. Para un canal Rayleigh con
desvanecimiento i.i.d., cada entrada de $\mathbf{H}$ es una variable compleja gaussiana de
media cero y varianza unitaria.

La capacidad ergódica, suponiendo que el transmisor no conoce el canal y reparte la
potencia por igual entre las $N_t$ antenas, es
\begin{equation}
  C = \mathbb{E}\left[\log_2\det\left(\mathbf{I}_{N_r}
      + \frac{\rho}{N_t}\mathbf{H}\mathbf{H}^{H}\right)\right],
\end{equation}
con $\rho$ la SNR lineal. Para SISO se reduce a
\begin{equation}
  C_{\mathrm{SISO}} = \mathbb{E}\left[\log_2\left(1+\rho|h|^2\right)\right].
\end{equation}
Si se descompone $\mathbf{H}=\mathbf{U}\boldsymbol{\Sigma}\mathbf{V}^{H}$, la capacidad se
escribe como
\begin{equation}
  C = \sum_{i=1}^{\min(N_t,N_r)}
      \log_2\left(1+\frac{\rho}{N_t}\sigma_i^2\right),
\end{equation}
es decir, cada valor singular abre un subcanal paralelo con ganancia $\sigma_i^2$. El
número de subcanales utilizables (ganancia de multiplexación) es
$\min(N_t,N_r)$.

\section{Desarrollo experimental}

\subsection{Parámetros}

\begin{tabularx}{\linewidth}{@{}*{2}{>{\raggedright\arraybackslash}X}@{}}
  \toprule
  Parámetro & Valor \\
  \midrule
  Configuraciones & SISO 1x1, MIMO 2x2, MIMO 4x4 \\
  SNR & 0, 5, 10, 15, 20, 25 y 30 dB \\
  Realizaciones por punto & 1000 \\
  Distribución del canal & Rayleigh i.i.d., $\mathcal{CN}(0,1)$ \\
  \bottomrule
\end{tabularx}

\subsection{Procedimiento}

\begin{enumerate}
  \item Para cada realización se genera $\mathbf{H}$ con entradas gaussianas complejas
        normalizadas.
  \item Se calcula la capacidad instantánea con la fórmula de determinante.
  \item Se promedian 1000 realizaciones para obtener la capacidad ergódica.
  \item Se obtienen los valores singulares de 1000 matrices 2x2 y 4x4 para estudiar su
        distribución.
\end{enumerate}

\section{Código}

\begin{codigo}
def canal(nr, nt, rng):
    return (rng.standard_normal((nr, nt))
            + 1j * rng.standard_normal((nr, nt))) / np.sqrt(2)

def capacidad(h, rho):
    nr, nt = h.shape
    return float(np.real(np.log2(np.linalg.det(
        np.eye(nr) + (rho / nt) * h @ h.conj().T))))
\end{codigo}

\section{Resultados}

\subsection{Capacidad versus SNR}

\begin{tabular}{@{}cccc@{}}
  \toprule
  SNR (dB) & SISO 1x1 & MIMO 2x2 & MIMO 4x4 \\
  \midrule
  0  & 0.86  & 1.69  & 3.35 \\
  5  & 1.77  & 3.28  & 6.53 \\
  10 & 2.96  & 5.56  & 10.89 \\
  15 & 4.31  & 8.28  & 16.25 \\
  20 & 5.81  & 11.37 & 22.17 \\
  25 & 7.57  & 14.51 & 28.32 \\
  30 & 9.18  & 17.78 & 34.98 \\
  \bottomrule
\end{tabular}

\figura{figs/capacidad.png}{Capacidad ergódica de canales Rayleigh para SISO, 2x2 y 4x4.}

La capacidad crece con la SNR y con el número de antenas. A 10 dB, el sistema 2x2 alcanza
5.56 bps/Hz frente a 2.96 del SISO (1.88 veces), y el 4x4 llega a 10.89 bps/Hz (3.68 veces
el SISO). Al duplicar el número de antenas de 2 a 4 la capacidad pasa de 5.56 a
10.89 bps/Hz, casi exactamente el doble.

\subsection{Valores singulares}

\figura{figs/valores_singulares.png}{Distribución de los valores singulares de canales 2x2 y 4x4.}

En el ejemplo 2x2 generado, los valores singulares fueron $\sigma_1=1.330$ y
$\sigma_2=0.684$. La capacidad instantánea a 10 dB para esa matriz es la suma de dos
términos, uno por cada subcanal. La dispersión de los valores singulares explica por qué la
capacidad instantánea varía tanto alrededor de la media: cuando los $\sigma_i$ son
similares, todos los subcanales aportan; cuando uno es muy pequeño, ese subcanal queda
casi inutilizable.

\section{Análisis de resultados}

\begin{itemize}
  \item \textbf{¿Por qué MIMO incrementa la capacidad?} Porque crea varios subcanales
        espaciales en la misma banda de frecuencia. Con $N_t$ antenas y $N_r$ antenas se
        pueden transmitir hasta $\min(N_t,N_r)$ flujos en paralelo, sin gastar más espectro
        ni más potencia total.
  \item \textbf{¿Cómo influye la SNR sobre el desempeño?} La capacidad crece de forma
        aproximadamente logarítmica con la SNR: cada 3 dB adicionales añaden del orden de
        1 bps/Hz por subcanal cuando la SNR es alta. A baja SNR los canales débiles aportan
        poco.
  \item \textbf{¿Qué sucede al duplicar el número de antenas?} La capacidad se
        aproximadamente duplica, porque se duplica el número de subcanales paralelos. En la
        simulación, 4x4 dio 34.98 bps/Hz frente a 17.78 de 2x2 a 30 dB.
  \item \textbf{¿Qué representan físicamente los valores singulares de la matriz de canal?}
        La ganancia (amplitud) de cada subcanal espacial independiente. Un valor singular
        grande indica una dirección en la que el canal transmite bien; uno pequeño, una
        dirección con desvanecimiento o alta atenuación.
\end{itemize}

\section{Cuestionario}

\begin{enumerate}
  \item \textbf{¿Qué significa MIMO?} Multiple Input Multiple Output: uso de múltiples
        antenas en transmisión y recepción para aumentar la capacidad, la confiabilidad o
        la cobertura de un enlace inalámbrico.
  \item \textbf{¿Cuál es la diferencia entre SISO y MIMO?} SISO usa una antena en cada
        extremo y un solo flujo; MIMO usa varias antenas y puede transmitir varios flujos
        simultáneos en la misma banda, además de aportar diversidad espacial.
  \item \textbf{¿Qué es la capacidad de canal?} La máxima tasa de información que se puede
        transmitir con una probabilidad de error arbitrariamente pequeña, dada una relación
        señal a ruido y un ancho de banda. Se mide en bits por segundo por hertzio.
  \item \textbf{¿Por qué LTE utiliza MIMO?} Para aumentar la tasa de datos y la robustez sin
        necesitar más espectro. LTE usa MIMO 2x2 y 4x4 en el enlace descendente, con
        multiplexación espacial y diversidad de transmisión.
  \item \textbf{¿Qué ventaja ofrece un sistema 4x4 respecto de un 2x2?} Aproximadamente el
        doble de capacidad en el mismo ancho de banda y con la misma SNR, porque ofrece
        cuatro subcanales espaciales en lugar de dos. El costo es mayor complejidad de
        procesamiento, más antenas y más consumo.
\end{enumerate}

\section{Conclusiones}

Se modelaron canales MIMO Rayleigh y se calculó su capacidad ergódica. A 10 dB se obtuvo
2.96 bps/Hz para SISO, 5.56 para 2x2 y 10.89 para 4x4; a 30 dB, 9.18, 17.78 y
34.98 bps/Hz, respectivamente. La capacidad crece de forma aproximadamente lineal con
$\min(N_t,N_r)$ y logarítmica con la SNR. Los valores singulares del canal explican la
variabilidad de la capacidad instantánea y confirman que MIMO aprovecha el espacio, no el
espectro.

\end{document}
